\chapter{A Hardware Primer for a DSP and ML Engineer}\label{ch:primer}

You already know the signal processing and the machine learning. This chapter
gives you the hardware concepts this project uses, each tied to a specific line
of your thesis code, with pointers to where to read more. Read it once quickly,
then come back to each section when a phase needs it.

% =========================================================================
\section{From algorithm to hardware: what changes}\label{sec:primer:change}

In Python, you pay for time, and everything else is nearly free. In hardware,
you pay for \emph{area} (silicon, or FPGA resources), \emph{energy} (battery
life) and \emph{design effort}, and you meet time as a hard deadline counted in
clock cycles. Four habits change:
\begin{itemize}
  \item \textbf{Numbers have a width you choose.} There is no free float64.
        Every bit costs area and energy (\cref{sec:primer:numbers}).
  \item \textbf{Operations are not equal.} An addition is cheap, a
        multiplication moderate, a division or \code{exp} expensive
        (\cref{sec:primer:arith}).
  \item \textbf{Memory costs more than arithmetic.} Fetching an operand from
        memory can cost far more energy than computing with it
        (\cref{sec:primer:memory}).
  \item \textbf{Everything happens in parallel.} Hardware is a circuit, not a
        sequence of statements. You decide how much of it to build and how to
        share it over time (\cref{sec:primer:arch}).
\end{itemize}

% =========================================================================
\section{Number representation}\label{sec:primer:numbers}

\begin{conceptbox}{Fixed point and Q-format}
A $W$-bit two's-complement integer $q$ represents the real value
$x = q\cdot2^{-F}$, where $F$ is the number of fractional bits. The notation
\textbf{Q$m$.$F$} means $m$ integer bits, including the sign, and $F$
fractional bits, with $W=m+F$. Signed Q$m$.$F$ covers $[-2^{m-1},\,2^{m-1}-2^{-F}]$
with resolution $2^{-F}$. The unsigned form UQ$m$.$F$ covers $[0,\,2^m-2^{-F}]$.
Examples: audio in Q1.15 covers $[-1,1)$; a gain in UQ3.13 covers $[0,8)$ with
resolution $1.2\times10^{-4}$.

When a result exceeds the range, \textbf{overflow} either \emph{wraps} (two's
complement arithmetic, which is catastrophic for audio) or \emph{saturates}
(clamps to the limit, which costs a comparator). When you discard low bits,
\textbf{rounding} either truncates (biased; free) or rounds to nearest (one
adder).
\end{conceptbox}

\begin{conceptbox}{Quantisation noise and SQNR}
Rounding to $F$ fractional bits adds an error that is roughly uniform on
$[-2^{-F-1},2^{-F-1}]$, with variance $\sigma^2=2^{-2F}/12$, provided the signal
is \enquote{busy} enough~\cite{widrow2008quantization}. For a full-scale sine
quantised to $N$ bits, this gives the familiar
\begin{equation}
  \text{SQNR} \approx 6.02\,N + 1.76\ \text{dB}.
  \label{eq:primer:sqnr}
\end{equation}
Each extra bit buys about \qty{6}{\dB}. The catch is \enquote{full-scale}:
signals that use only part of the range waste bits. That is why the dynamic
range profiling in \cref{ch:codesign} comes first.
\end{conceptbox}

\begin{whybox}[Where this appears in your thesis code]
The Wiener gain at \file{wiener.py:102} lies between the floor (0.05) and 1
before tilt and output gain, and up to 4 after. In UQ3.13 its quantisation step
is $1.2\times10^{-4}$, about \qty{0.001}{\dB} near unity gain but about
\qty{0.02}{\dB} at the 0.05 floor, because the step is fixed in absolute terms.
A logarithmic format would keep the relative step constant. Questions like this
are exactly what Phase~1 measures.
\end{whybox}

\begin{conceptbox}{Block floating point, small floats and logarithmic numbers}
\begin{itemize}
  \item \textbf{Block floating point}: a group of values, such as an FFT frame,
        shares one exponent. This is cheap and very effective in FFTs.
  \item \textbf{Small custom floats}, for example a 10-bit mantissa and a 6-bit
        exponent: the relative precision is constant over a huge range, which is
        good for power spectra.
  \item \textbf{Logarithmic number systems}: store $\log_2|x|$. Multiplication
        becomes addition and division becomes subtraction, while addition
        becomes harder. Comparisons and minima are exact, since $\log$ is
        monotone. Your IMCRA minimum trackers are a natural fit.
\end{itemize}
\end{conceptbox}

% =========================================================================
\section{Arithmetic building blocks and their cost}\label{sec:primer:arith}

\begin{table}[htbp]
\centering\footnotesize
\caption{Relative cost of operations in hardware (qualitative; it depends on
width and technology).}
\label{tab:primer:cost}
\begin{tabularx}{\textwidth}{@{}lLL@{}}
\toprule
Operation & Typical implementation & Cost \\
\midrule
Add, subtract, compare & ripple or carry-lookahead adder & low: linear in width \\
Shift by a constant & wiring & free \\
Multiply & array or Booth multiplier; a DSP slice on an FPGA & moderate: roughly quadratic in width \\
Divide & iterative (one bit per cycle), or reciprocal + multiply & high, or slow \\
$\exp$, $\log$, $x^\beta$, $\sqrt{\cdot}$ & table + interpolation, CORDIC, or polynomial & high; depends on accuracy \\
Table lookup & ROM or RAM & cheap if small; cost grows with size \\
\bottomrule
\end{tabularx}
\end{table}

\begin{conceptbox}{Reciprocal by Newton--Raphson}
To compute $y=1/d$, apply Newton's method to $f(y)=1/y-d$:
\begin{equation}
  y_{k+1} = y_k\,(2 - d\,y_k).
\end{equation}
The relative error squares at every step: if $y_k=(1-e)/d$ then
$y_{k+1}=(1-e^2)/d$. So a seed from a small table with 8 correct bits becomes 16
after one step, at the cost of two multiplications. Normalise $d$ first with a
leading-zero count, so the table only has to cover $[1,2)$, then shift the
result back~\cite{muller2016elementary}.
\end{conceptbox}

\begin{conceptbox}{Tables, interpolation and CORDIC}
A smooth function on a bounded range can be stored as a table of $2^a$ samples.
Linear interpolation between neighbouring entries costs one multiply and greatly
improves accuracy for the same table size. $\log_2$ and $2^x$ are especially
cheap: the integer part comes from the leading-zero count (a priority encoder)
and only the fractional part needs a table. \emph{CORDIC} computes trigonometric,
hyperbolic, exponential and logarithmic functions with only shifts and
additions, one bit of accuracy per iteration~\cite{volder1959cordic}. It is
compact, but needs many cycles.
\end{conceptbox}

\begin{whybox}[Where this appears in your thesis code]
Per bin and per frame, the thesis computes 12 divisions, one \code{exp}
(\file{imcra.py:153}) and one \code{pow} (\file{wiener.py:102})
(\cref{ch:workload}). Phase~1 replaces them with shared reciprocals, small
tables and a gain lookup table rebuilt by the control plane.
\end{whybox}

% =========================================================================
\section{The FFT in hardware}\label{sec:primer:fft}

\begin{conceptbox}{The radix-2 butterfly}
The Cooley--Tukey FFT~\cite{cooley1965fft} splits an $N$-point DFT into two
$N/2$-point DFTs, recursively, joined by \emph{butterflies}:
$A' = A + W_N^k B$ and $B' = A - W_N^k B$. An $N$-point transform has $\log_2 N$
stages of $N/2$ butterflies. Each butterfly is one complex multiplication (four
real multiplications and two additions) and two complex additions. Hardware
needs the \emph{twiddle factors} $W_N^k$, which are stored in a ROM; by symmetry
a quarter wave is enough. It also needs \emph{bit-reversed} ordering at the input
or output.
\end{conceptbox}

There are two families of FFT architecture:
\begin{itemize}
  \item \textbf{Memory-based (in place).} One or a few butterfly units read and
        write a RAM over $\log_2 N$ passes. It is small and slow, which is ideal
        when you have many cycles per frame, as you do.
  \item \textbf{Pipelined}, such as single-path delay feedback (R2SDF,
        R2\textsuperscript{2}SDF). There is one butterfly per stage, connected by
        delay lines, and the unit accepts one sample per cycle
        continuously~\cite{he1996r22sdf,ch12-13:garrido2022fft}. It is larger,
        and the choice for streaming at high sample rates.
\end{itemize}
Values can grow by one bit per stage. Either provide $\log_2 N$ guard bits,
scale by $1/2$ at every stage (which loses precision for small signals), or use
block floating point, scaling only when growth is detected. A \textbf{real-input
FFT} of size $N$ can be computed with an $N/2$-point complex FFT and a cheap
post-processing step, which halves the work for audio~\cite{oppenheim2010dtsp}.

% =========================================================================
\section{Architecture transformations}\label{sec:primer:arch}

\begin{conceptbox}{Throughput, latency and the initiation interval}
\emph{Latency} is how many cycles one input takes to produce its output.
\emph{Throughput} is how many inputs are accepted per cycle. A design with
initiation interval II accepts a new input every II cycles. Pipelining
(inserting registers) shortens the critical path, allowing a faster clock, and
keeps II at 1, at the cost of more latency and registers.
\end{conceptbox}

\begin{conceptbox}{Folding and unfolding}
\emph{Folding} (resource sharing) maps several operations of the algorithm onto
one hardware unit over several cycles: less area, more cycles. \emph{Unfolding}
(parallelism) does the opposite. \emph{Retiming} moves registers to balance
paths without changing function. These are the main tools of VLSI DSP
architecture~\cite{parhi1999vlsidsp}.
\end{conceptbox}

\begin{whybox}[A worked example from your thesis]
The per-bin processing is about 112 operations for each of 1025 bins
(\cref{ch:workload}).
\begin{itemize}
  \item A \emph{fully pipelined} unit (II = 1, all operations in hardware)
        processes the bins in about \num{1025} cycles plus its pipeline depth.
  \item A \emph{fully folded} unit, one arithmetic unit doing about one
        operation per cycle, needs about \num{115000} cycles.
\end{itemize}
At \qty{10}{\MHz} the frame budget is \num{106667} cycles, so the folded design
just misses it, while the pipelined one uses about 1\%. Between these extremes
lies a whole family of designs. Finding the one with the least energy per frame
is part of RQ3.
\end{whybox}

% =========================================================================
\section{Memories}\label{sec:primer:memory}

\begin{itemize}
  \item \textbf{Flip-flops} hold one bit each, can be read in parallel, and are
        expensive per bit. Use them for pipelines and small state.
  \item \textbf{SRAM} is dense, with one or two ports, each reading or writing
        one word per cycle. On an FPGA it appears as block RAM, which on the
        Zynq-7020 comes in 36\,Kbit blocks, or as distributed \enquote{LUT RAM}.
  \item \textbf{ROM} holds constant tables. On FPGAs it is block RAM or logic;
        in ASICs, logic or a ROM macro.
\end{itemize}
Accessing memory typically costs considerably more energy than an arithmetic
operation on the same data, and the gap grows with memory
size~\cite{ch12-13:horowitz2014energy}. Two rules follow:
\begin{itemize}
  \item make memories small, which is why word lengths and bin counts matter;
  \item touch them rarely, which is why caching the tracker minima matters
        (\cref{ch:workload}).
\end{itemize}

% =========================================================================
\section{Power and energy}\label{sec:primer:power}

\begin{conceptbox}{Dynamic and static power}
CMOS logic consumes dynamic power each time a node switches:
\begin{equation}
  P_\text{dyn} = \alpha\,C\,V_{DD}^2\,f,
\end{equation}
with activity factor $\alpha$, switched capacitance $C$, supply voltage
$V_{DD}$ and clock frequency $f$. \emph{Leakage} (static) power flows whenever
the circuit is powered, and is ever more important in small processes.
\end{conceptbox}

Since the work per frame is fixed, the right metric is \textbf{energy per
frame}, not power. Running faster does not save dynamic energy, since the same
nodes switch either way. But a design that can run \emph{slower} can run at a
\emph{lower voltage}, and energy scales with $V^2$. That is why low-power DSP
uses parallel or pipelined hardware at a low clock and a low
voltage~\cite{chandrakasan1992lowpower}. Clock gating stops switching in idle
blocks, and power gating removes leakage from them. Your datapath is idle for
most of each frame, so both matter.

% =========================================================================
\section{Implementation targets}\label{sec:primer:targets}

\begin{table}[htbp]
\centering\footnotesize
\caption{Where a signal-processing algorithm can run.}
\label{tab:primer:targets}
\begin{tabularx}{\textwidth}{@{}lLLL@{}}
\toprule
Target & Strength & Weakness & In this project \\
\midrule
DSP processor or MCU & Flexible software; fast to develop & Energy per operation well above dedicated hardware & The comparison point; the control plane on the ARM \\
MCU + small NPU & Neural networks at low power & Fixed set of supported operators & Context for Track A \\
FPGA & Reconfigurable; real-time prototyping & Larger power and area than an ASIC & Phase~4 prototype \\
ASIC & Lowest energy and area & Expensive, slow, unforgiving & Track B, with an open PDK \\
\bottomrule
\end{tabularx}
\end{table}

% =========================================================================
\section{Design entry: RTL, HLS or a generator}\label{sec:primer:entry}

\begin{itemize}
  \item \textbf{RTL in SystemVerilog}~\cite{sutherland2006svdesign}: you
        describe registers and logic explicitly. This is the industry standard
        and what VLSI and ASIC roles expect you to read and write.
  \item \textbf{High-level synthesis (HLS)}: C++ compiled to hardware. It is
        fast for dataflow and ML kernels (hls4ml~\cite{duarte2018hls4ml} builds
        on it), but harder to control at the bit and cycle level.
  \item \textbf{Hardware generators} such as Chisel, Amaranth or SpinalHDL:
        hardware described in Scala or Python. They are powerful, but less
        common in DSP-ASIC teams.
\end{itemize}

\textbf{Recommendation:} hand-written SystemVerilog for the DSP datapath, so
you learn the craft that roles ask for. HLS is optional for Track A.

% =========================================================================
\section{Interfaces}\label{sec:primer:interfaces}

\begin{itemize}
  \item \textbf{Valid/ready handshake.} A transfer happens on a cycle when both
        \sig{valid} (the producer has data) and \sig{ready} (the consumer can
        take it) are high. It is the universal streaming interface inside
        designs.
  \item \textbf{AXI4-Stream}: valid/ready plus data and optional side-band
        signals, used with DMA engines on Zynq.
  \item \textbf{AXI4-Lite}: a simple memory-mapped register interface, with
        separate address and data channels for writes and reads. It is how the
        ARM writes your registers (\cref{app:registers}).
  \item \textbf{I2S}: a serial audio interface with a bit clock, a word-select
        (left/right) clock and a data line, used between the codec and the FPGA.
\end{itemize}

% =========================================================================
\section{Verification basics}\label{sec:primer:verif}

A \emph{testbench} drives a design and checks its outputs against a
\emph{reference model}. When the model is bit-accurate, the check is exact
equality. \emph{Code coverage} measures which lines and toggles were exercised.
\emph{Functional coverage} measures which situations you intended to test
actually occurred. \emph{Assertions} state invariants that simulation checks
continuously. \emph{Formal verification} proves them for all possible inputs,
which is practical for small blocks. \Cref{ch:verification} puts all of this
into practice with cocotb~\cite{cocotb}, Verilator~\cite{verilator} and
SymbiYosys.

% =========================================================================
\section{ML inference hardware basics}\label{sec:primer:ml}

\begin{itemize}
  \item \textbf{MAC arrays.} Neural-network inference is dominated by
        multiply--accumulate operations. Accelerators are arrays of MAC units
        fed from local memories.
  \item \textbf{Systolic arrays.} A grid of MAC cells passes data to its
        neighbours every cycle, so each value fetched from memory is reused many
        times~\cite{kung1982systolic}. A TPU is built around
        one~\cite{jouppi2017tpu}.
  \item \textbf{Dataflows.} \emph{Weight-stationary} keeps weights in the cells;
        \emph{output-stationary} keeps partial sums; \emph{row-stationary}
        balances both~\cite{chen2017eyeriss}. The choice decides which memory
        traffic you save.
  \item \textbf{Quantisation.} Int8 weights and activations with int32
        accumulation are standard, with a small loss in accuracy for most
        models~\cite{jacob2018quantization}.
  \item \textbf{Roofline.} Achievable performance is the minimum of peak
        compute and memory bandwidth times arithmetic intensity, the operations
        per byte moved~\cite{williams2009roofline}. Batch-1 matrix--vector
        products have low intensity, so they are memory-bound.
\end{itemize}

\begin{readbox}[Read --- the essential four]
Harris and Harris, for digital design~\cite{harris2021ddca}. Parhi, for VLSI
DSP architecture~\cite{parhi1999vlsidsp}. Sze \etal, for ML
hardware~\cite{sze2017efficient}. Widrow and Koll\'ar, for
quantisation~\cite{widrow2008quantization}. \Cref{app:reading} has the full
annotated list.
\end{readbox}

\begin{donebox}[You have what you need from this chapter when you can]
\begin{checklist}
  \item read Q3.13 and compute its range and resolution in your head;
  \item explain why one Newton--Raphson step doubles the correct bits;
  \item sketch a radix-2 butterfly and say how many there are in a 1024-point FFT;
  \item explain folding versus pipelining with the bin-processing example;
  \item explain why energy per frame, not power, is the right metric here.
\end{checklist}
\end{donebox}
